\section{DATA PREPARATION}
The datasets used in this project are MedQA-USMLE \cite{jin2021} and MedMCQA \cite{pal2022}, both of which are also used for benchmark evaluation in the AMG-RAG methodology of \cite{rezaei2025}. MedQA is a multiple-choice question answering dataset derived from professional medical licensing examinations, and its questions were written in reference to a collection of 18 standard medical textbooks. MedMCQA is built from medical entrance examination data and covers foundational and clinical knowledge. For the purposes of this system, the MedQA textbooks serve as the RAG retrieval corpus, while MedMCQA is used exclusively to train the specialty classifier.

Initial exploratory analysis confirmed answer balance in MedQA and characterised the specialty distribution in MedMCQA. Textbook character counts and chunk distributions were subsequently verified prior to vectorstore construction.

\begin{figure}[htbp]
    \centering
    \includegraphics[width=\textwidth]{textbooks.png}
    \caption{Left: Character count distribution across the 18 MedQA textbooks. Right: Top 15 subjects of MedMCQA answers.}
    \label{fig:textbooks}
\end{figure}

Following exploratory checks, the MedQA textbooks are chunked into 400-word sections with a 50-word overlap, producing 36,723 chunks. To support the embedding model ablation described in Section 8, three separate FAISS vectorstores were constructed: one each for all-MiniLM-L12-v2, Qwen3-Embedding-0.6B, and Octen-Embedding-0.6B.